\documentclass[11pt,a4paper]{article}
\usepackage[T1]{fontenc}
\usepackage{lmodern,amsmath,amssymb,geometry}
\usepackage[hidelinks]{hyperref}
\geometry{margin=2.5cm}
\setlength{\parindent}{0pt}
\setlength{\parskip}{0.6em}
\begin{document}
\section*{Explicit midpoint}
For $z'=f(t,z)$ the two stages and accepted state are
\[
k_1=f(t_n,z_n),\qquad k_2=f(t_n+h/2,z_n+hk_1/2),\qquad z_{n+1}=z_n+hk_2.
\]
Thus $A_{21}=1/2$, all other entries of $A$ vanish, $b=(0,1)$ and $c=(0,1/2)$.
The order-two conditions $\sum_i b_i=1$ and $\sum_i b_ic_i=1/2$ hold.
For sufficiently smooth $f$, the local truncation error is $O(h^3)$ and global error is $O(h^2)$ on bounded time intervals under the usual stability assumptions.
Each step uses two field evaluations and no algebraic solve.

For the harmonic oscillator $z'=Jz$, $J^2=-I$, the map is
$M_h=(1-h^2/2)I+hJ$ and $\det M_h=1+h^4/4$.
Hence this explicit midpoint method is not symplectic and does not preserve the oscillator's quadratic energy for nonzero $h$.

\section*{Implementation boundary}
The named recipe inherits the shared Runge--Kutta stage algorithm and declares only its fixed tableau and order.
Methods receive a field and independent initial conditions; physical dynamics and passive energy observation remain outside the step.
Execution owns fixed-step scheduling, output sampling and NumPy/JAX selection.
See the adjacent simulation guide for public construction and configuration.
\end{document}
